% Bagian: first-order-logic
% Bab: beyond
% Subbagian: higher-order-logic

\documentclass[../../../include/open-logic-section]{subfiles}

\begin{document}

\olfileid[id]{fol}{byd}{hol}

\olsection{Logika Orde Tinggi}

Peralihan dari logika orde pertama ke logika orde kedua memungkinkan kita
membicarakan himpunan objek dalam !!{domain} orde pertama di dalam bahasa
formal. Mengapa berhenti di sana? Misalnya, logika orde ketiga semestinya
memungkinkan kita menangani himpunan dari himpunan objek, atau bahkan himpunan
yang memuat baik objek maupun himpunan objek. Dan logika orde keempat akan
memungkinkan kita membicarakan himpunan objek semacam itu. Seperti mungkin
telah Anda duga, gagasan ini dapat diiterasikan tanpa batas.

Dalam praktiknya, logika orde tinggi sering di!!{formula}{}sikan dalam istilah
fungsi, bukan relasi. (Dengan memperhitungkan identifikasi alaminya, perbedaan
ini tidak hakiki.) Diberikan beberapa ``sorta'' dasar $A$, $B$, $C$,~\dots
(yang sekarang akan kita sebut ``tipe''), kita dapat membuat tipe baru dengan
menetapkan
\begin{quote}
Jika $\sigma$ dan $\tau$ merupakan tipe hingga, maka $\sigma \to \tau$ juga
merupakan tipe hingga.
\end{quote}
Bayangkan tipe sebagai ``label'' sintaktis yang mengklasifikasikan objek yang
kita inginkan dalam !!{domain}; $\sigma \to \tau$ menggambarkan objek-objek
yang merupakan fungsi yang membawa objek bertipe~$\sigma$ ke objek
bertipe~$\tau$. Misalnya, kita mungkin ingin mempunyai tipe $\Omega$ untuk
nilai kebenaran, ``benar'' dan ``salah'', serta tipe $\Nat$ untuk bilangan
asli. Dalam hal itu, objek bertipe $\Nat \to \Omega$ dapat dipandang sebagai
relasi uner, atau subhimpunan dari $\Nat$; objek bertipe $\Nat \to \Nat$
merupakan fungsi dari bilangan asli ke bilangan asli; dan objek bertipe
$(\Nat \to \Nat) \to \Nat$ merupakan ``fungsional'', yakni fungsi bertipe
tinggi yang membawa fungsi ke bilangan.

Seperti dalam logika orde kedua, logika orde tinggi dapat dipandang sebagai
sejenis logika banyak-sorta, dengan satu sorta untuk setiap tipe objek yang
ingin kita pertimbangkan. Namun, biasanya lebih jelas untuk langsung
mendefinisikan sintaks logika bertipe tinggi dari dasar. Misalnya, kita dapat
mendefinisikan suatu himpunan tipe hingga secara induktif sebagai berikut:
\begin{enumerate}
\item $\Nat$ merupakan tipe hingga.
\item Jika $\sigma$ dan $\tau$ merupakan tipe hingga, maka $\sigma
  \to \tau$ juga merupakan tipe hingga.
\item Jika $\sigma$ dan $\tau$ merupakan tipe hingga, maka $\sigma \times
  \tau$ juga merupakan tipe hingga.
\end{enumerate}
Secara intuitif, $\Nat$ menyatakan tipe bilangan asli, $\sigma \to \tau$
menyatakan tipe fungsi dari $\sigma$ ke $\tau$, dan $\sigma \times \tau$
menyatakan tipe pasangan objek, yang satu dari $\sigma$ dan yang lain dari
$\tau$. Selanjutnya, kita dapat mendefinisikan suatu himpunan suku secara
induktif sebagai berikut:
\begin{enumerate}
\item Untuk setiap tipe $\sigma$, tersedia sekumpulan variabel $x$, $y$,
  $z$, \dots bertipe $\sigma$.
\item $\Obj 0$ merupakan suku bertipe $\Nat$.
\item $\Obj S$ (penerus) merupakan suku bertipe $\Nat \to \Nat$.
\item Jika $s$ merupakan suku bertipe $\sigma$, dan $t$ merupakan suku
  bertipe $\Nat \to (\sigma \to \sigma)$, maka $\Obj{R}_{st}$ merupakan
  suku bertipe $\Nat \to \sigma$.
\item Jika $s$ merupakan suku bertipe $\tau \to \sigma$ dan $t$ merupakan
  suku bertipe~$\tau$, maka $s(t)$ merupakan suku bertipe $\sigma$.
\item Jika $s$ merupakan suku bertipe~$\sigma$ dan $x$ merupakan variabel
  bertipe~$\tau$, maka $\lambd[x][s]$ merupakan suku bertipe $\tau \to
  \sigma$.
\item Jika $s$ merupakan suku bertipe~$\sigma$ dan $t$ merupakan suku
  bertipe~$\tau$, maka $\tuple{s, t}$ merupakan suku bertipe $\sigma \times
  \tau$.
\item Jika $s$ merupakan suku bertipe~$\sigma \times \tau$, maka $p_1(s)$
  merupakan suku bertipe~$\sigma$ dan $p_2(s)$ merupakan suku
  bertipe~$\tau$.
\end{enumerate}
Secara intuitif, $\Obj{R}_{st}$ menyatakan fungsi yang didefinisikan secara
rekursif oleh
\begin{align*}
\Obj{R}_{st}(0) & = s \\
\Obj{R}_{st}(x+1) & = t(x, R_{st}(x)),
\end{align*}
$\tuple{s, t}$ menyatakan pasangan yang komponen pertamanya~$s$ dan komponen
keduanya~$t$, sedangkan $p_1(s)$ dan~$p_2(s)$ menyatakan elemen pertama dan
kedua (``proyeksi'') dari~$s$. Terakhir, $\lambd[x][s]$ menyatakan fungsi~$f$
yang didefinisikan oleh
\[
f(x) = s
\]
untuk setiap~$x$ bertipe~$\tau$; dengan demikian, butir~(6) memberikan suatu
bentuk komprehensi yang memungkinkan kita mendefinisikan fungsi menggunakan
suku. !!^{formula}s dibangun dari pernyataan !!{identity} $\eq[s][t]$ antara
suku-suku bertipe sama, penghubung proposisional biasa, dan kuantifikasi tipe
tinggi. Aksioma sistem kemudian dapat diambil sebagai persamaan dasar yang
mengatur suku-suku yang didefinisikan di atas, bersama aturan logika biasa
dengan kuantor dan !!{identity}.

Jika sistem tipe hingga diperluas dengan tipe nilai kebenaran~$\Omega$, kita
harus menyertakan pula aksioma-aksioma yang mengatur penggunaannya. Bahkan,
jika cukup cermat, kita dapat menghapus !!{formula}s kompleks sama sekali dan
menggantinya dengan suku bertipe~$\Omega$!{} Sistem bukti kemudian dapat
diubah sebagaimana mestinya. Hasilnya pada hakikatnya adalah \emph{teori tipe
sederhana} yang dikemukakan Alonzo Church pada tahun 1930-an.

Seperti dalam logika orde kedua, terdapat berbagai versi semantik tipe tinggi
yang mungkin ingin digunakan. Dalam versi penuh, variabel bertipe $\sigma \to
\tau$ merentang atas himpunan \emph{semua} fungsi dari objek bertipe~$\sigma$
ke objek bertipe~$\tau$. Seperti yang mungkin diduga, semantik ini terlalu
kuat untuk mempunyai sistem !!{derivation} efektif yang lengkap. Namun, kita
dapat mempertimbangkan semantik yang lebih lemah, dengan !!a{structure}
terdiri atas himpunan elemen $T_\tau$ bagi setiap tipe~$\tau$, bersama operasi
yang sesuai bagi aplikasi, proyeksi, dan seterusnya. Jika semua perinciannya
dikerjakan dengan benar, kita dapat memperoleh teorema kelengkapan bagi
jenis-jenis sistem !!{derivation} yang diuraikan di atas.

Logika bertipe tinggi menarik karena menyediakan kerangka yang memungkinkan
kita menanamkan banyak bagian matematika secara alami: dengan berawal dari
$\Nat$, kita dapat mendefinisikan bilangan real, fungsi kontinu, dan
seterusnya. Logika ini juga khususnya menarik dalam konteks logika
intuisionistik, karena tipe mempunyai interpretasi ``konstruktif'' yang jelas.
Bahkan, kita dapat mengembangkan versi konstruktif semantik tipe tinggi
(berdasarkan logika intuisionistik, bukan logika klasik) yang memperjelas
interpretasi konstruktif ini dengan sangat baik dan, dalam banyak hal, lebih
menarik daripada padanannya yang klasik.

\end{document}
